\section{Exact composition with the literal receiver}\label{sec:receiver-bridge}

The trace return and the receiver bounds have different domains.  This
section gives their exact composite and records the data that the composite
retains or forgets.

Let \(\mathcal T_{24}^{\mathrm{mix}}\) be the set of marked tuples
\[
 T=(p,c,q,a,y,z,\omega)
\]
such that
\[
 p\equiv1\pmod {24},\quad a=cq,\quad c\mid6,\quad q>3\text{ prime},
 \quad \frac p4<a<\frac p2,
\]
\[
 y,z\in\mathbb Q_{>0},\qquad
 \frac4p=\frac1a+\frac1y+\frac1z,qquad y+z\in\mathbb Z,
\]
where \(\omega\) records the original order of the two tails.  The
factorization in Lemma~\ref{lem:trace} assigns an E or M word
\[
 R=4a-p,qquad u\mid a^2,qquad
 a=hrs,qquad u=hr^2,qquad \gcd(r,s)=1.
\]
Its common reduced tail denominator is
\[
 d(T)=\frac{R}{\gcd(R,G_E)}
 \quad\hbox{or}\quad
 d(T)=\frac{R}{\gcd(R,G_M)},
\]
according to the channel.  Let \(\mathcal W_{12}^{\mathrm{ord}}\) be the
ordered integral witnesses
\[
 W=(p;a_W,Y_W,Z_W),\qquad p\equiv1\pmod {12},qquad
 \frac4p=\frac1{a_W}+\frac1{Y_W}+\frac1{Z_W},qquad
 a_W<Y_W<Z_W.
\]
Then, after matching the retained order,
\[
 \mathcal T_{24}^{\mathrm{mix}}\cap\mathcal W_{12}^{\mathrm{ord}}
 =\{T\in\mathcal T_{24}^{\mathrm{mix}}:d(T)=1\}.
\]
In particular, a proper trace source \(d(T)>1\) is not already in the
receiver theorem's integral domain.

The mixed theorem defines a source-changing map
\begin{equation}\label{eq:rho}
 \rho:\mathcal T_{24}^{\mathrm{mix}}
 \longrightarrow \mathcal W_{24}^{M,\mathrm{ord}}
 \subset\mathcal W_{12}^{\mathrm{ord}}.
\end{equation}
For \(W=\rho(T)\), retain
\[
 \ell=(p,a_W,Y_W,Z_W),\qquad S=p+a_W+Y_W+Z_W,qquad
 H_W(X)=-S^{-1}\prod_{j=0}^3(X-\ell_j).
\]
Write \(d_j=H_W'(\ell_j)\), and retain one of the two square roots
\(\xi_j\) with \(\xi_j^2d_j=1\) for every label.  The literal receiver is
the typed linear map
\begin{equation}\label{eq:receiver-typed}
 \mathscr O(W,\xi):V_\beta=\mathbb C^4_{\mathrm{labels}}
 \longrightarrow E_q=\mathbb C^4_q,
\end{equation}
\[
 e_j\longmapsto
 \xi_j
 \begin{pmatrix}
 1\\-i d_j\\A d_j^2+2\ell_jd_j\\
 i(7\ell_j^2d_j-13d_j^2)
 \end{pmatrix},
 \qquad A=-S^{-1}.
\]
Let \(\Xi_W\) be the sixteen retained square-root choices.  The exact
composite is therefore
\begin{equation}\label{eq:receiver-composite}
 \mathscr O\circ\rho:
 \mathcal T_{24}^{\mathrm{mix}}\times_\rho\Xi
 \longrightarrow
 \operatorname{Iso}_{\mathbb C}(V_\beta,E_q).
\end{equation}
The codomain follows from the global-sign theorem and the sharp inverse
theorem~\cite{GlobalSign,SharpGrowth}.  It is not assumed in defining
\(\rho\).

\subsection{The two branches in receiver coordinates}

For an E input, write
\[
 m=4rs,qquad R=\delta t^2,qquad
 H=\frac{\delta+1}{m},qquad J=\frac{pr+s}{\delta}.
\]
The returned M state before sorting has denominators
\[
 (HsJ,\ pHJr,\ pHsr).
\]
Since \(J\ge r+s>s\), the receiver order is
\begin{equation}\label{eq:E-return-sorted}
 W_M=(p;HsJ,pHsr,pHJr).
\end{equation}
Its literal receiver parameters are
\begin{equation}\label{eq:E-return-parameters}
 R_M=\frac{s+J}{r},\qquad
 w_M=\frac{p}{R_M}=\frac{pr}{s+J},qquad
 b=Hsr,qquad c=HJr,qquad
 \tau=\frac{pHr(J-s)}{S}.
\end{equation}
The squarefree number
\[
 \delta=4Hsr-1
\]
is the retained cofactor \(Q_M\), not the receiver residual \(R_M\).
Replacing \(w_M\) by \(p/\delta\) would change the receiver theorem.  If
the input order had the exterior-swap bit, the target complement
\(a_M^2/u_M=HJ^2\) restores that order before the independent sorting in
\eqref{eq:E-return-sorted}.

For an M input, let \(u_0\mid c^2\mid36\) be the small word after the marked
complement, and put
\[
 a'=\frac{p+\delta}{4},\qquad g'=\gcd(a',u_0),qquad
 (h',r',s')=\left(\frac{g'^2}{u_0},\frac{u_0}{g'},\frac{a'}{g'}\right),
 \qquad \lambda'=\frac{r'+s'}{\delta}.
\]
Before the complement bit is restored, the return is
\[
 (a',ph's'\lambda',ph'r'\lambda').
\]
The receiver order is
\begin{equation}\label{eq:M-return-sorted}
 Y_W=ph'\lambda'\min(r',s'),qquad
 Z_W=ph'\lambda'\max(r',s'),qquad w=\frac p\delta.
\end{equation}
The M gate used here is the same gate in both presentations:
\[
 R\mid p+4u\quad\Longleftrightarrow\quad R\mid a+u,
\]
because \(p+4u=4(a+u)-R\) and \(R\) is odd.

\subsection{Quantitative theorem on the returned witness}

Put
\[
 c_{24}=\frac{327448292668}{47045881},
 \qquad \gamma=\frac{80}{6279^4}.
\]
For every \(T\in\mathcal T_{24}^{\mathrm{mix}}\), the integral witness
\(W=\rho(T)\) satisfies
\begin{equation}\label{eq:bridge-global-sign}
 \mathfrak N(W)<-c_{24}p^8,qquad
 |\det\mathscr O_W|>\frac{4c_{24}p^8}{S^3}.
\end{equation}
Because \(\rho(T)\) is a middle-channel witness, the middle-channel proof of
the sharp theorem gives
\begin{align}
 \|\mathscr O_W^{-1}\|_2
 &<\frac{224}{\gamma}p\sqrt w,
 &\|\wedge^2\mathscr O_W^{-1}\|_2
 &<\frac{180}{\gamma},\label{eq:bridge-inverse}\tag{RB1}\\
 \|\wedge^3\mathscr O_W^{-1}\|_2
 &<\frac{41}{4\gamma p^2},
 &|\det\mathscr O_W^{-1}|
 &<\frac1{4\gamma p^2S^3}.label{eq:bridge-compounds}\tag{RB2}
\end{align}
Here \(w\) and \(\tau\) are the actual returned-witness parameters in
\eqref{eq:E-return-parameters} or \eqref{eq:M-return-sorted}.  The two
determinant estimates combine without discarding either shape term:
\begin{equation}\label{eq:bridge-det-max}
 |\det\mathscr O_W|>
 \max\left\{
  \frac{4c_{24}p^8}{S^3},
  4\gamma p^2S^3\max(1,w^2\tau^2)
 \right\}.
\end{equation}
Moreover the third singular value obeys
\[
 s_3(\mathscr O_W)>\sqrt{\gamma/180},
\]
so at most one inverse singular direction can become unbounded.

\subsection{Fibres, symmetries, and exact limits of the composition}

For fixed \(p,r,s,\delta\), all E trace sources
\[
 R=\delta t^2,qquad t\mid A/\delta,qquad \delta t^2<p
\]
have the same returned E and M witnesses.  On the free modules of
Section~\ref{sec:auto}, the return therefore factors through \(\pi\), whose
kernel has basis
\[
 e_{\delta,t}-e_{\delta,1}\qquad(t>1).
\]
Retaining \(t\) before applying \(\pi\) recovers the original residual and
shell.  The E-to-E return preserves \(p,r,s\) and every centred exponent
\(v_\ell(u)-v_\ell(a)\).  The E-to-M return instead encodes
\[
 r=\lambda_M,qquad s=r_M,qquad J=s_M,qquad
 \delta=4h_Mr_M\lambda_M-1.
\]
The M return retains \(p,u_0\) and the complement bit, while recomputing the
target exponent box at \(a'\).

Tail exchange acts on \(\mathscr O\) by a column permutation.  Changing any
square-root branch acts by a diagonal sign matrix on the label side.  Hence
\(\mathfrak N\), the singular values, and the norm bounds are unchanged;
the oriented sign \(\epsilon\) and \(\det\mathscr O\) change compatibly.
The prime-indexed centred exponent vector is not identified with the
four-dimensional label vector \(V_\beta\); no map making that identification
has been proved.

The composition begins with a trace source.  It does not prove that
\(\mathcal T_{24}^{\mathrm{mix}}(p)\) is nonempty.  Receiver invertibility
does not imply a trace gate, a full gate, a ray deletion, or universal
occupancy.  The sharpness families for the receiver exponents lie in
\(p\equiv13\pmod{24}\), so their optimality does not prove optimality on the
restricted image of \(\rho\).

\begin{figure}[ht]
\centering
\resizebox{\textwidth}{!}{%
\begin{tikzpicture}[
  node distance=13mm and 16mm,
  >=Latex,
  box/.style={draw,rounded corners,align=center,inner sep=6pt,text width=37mm},
  map/.style={->,thick},
  note/.style={align=center,font=\small,text width=41mm}
]
\node[box] (trace) {marked rational trace\\
\(T=(p,c,q,a,y,z,\omega)\)\\
\(y+z\in\mathbb Z\)};
\node[box,right=of trace] (factor) {exact factorization\\
\(R,u,h,r,s,d\)\\
E or M channel};
\node[box,above right=9mm and 16mm of factor] (eray) {E branch\\
\(R=\delta t^2,\ rs\mid6\)\\
\(H=(\delta+1)/(4rs)\), \(J=(pr+s)/\delta\)};
\node[box,below right=9mm and 16mm of factor] (mword) {M branch\\
marked complement\\
\(u_0\mid c^2\mid36\), \(a'=(p+\delta)/4\)};
\node[box,right=22mm of factor] (witness) {integral M witness \(W=\rho(T)\)\\
E input: \((HsJ,pHsr,pHJr)\)\\
M input: \eqref{eq:M-return-sorted}};
\node[box,right=of witness] (receiver) {literal receiver\\
\(\mathscr O(W,\xi):V_\beta\to E_q\)\\
invertible with \eqref{eq:bridge-inverse}--\eqref{eq:bridge-det-max}};

\draw[map] (trace) -- node[above,font=\small] {Lemma~\ref{lem:trace}} (factor);
\draw[map] (factor) -- (eray);
\draw[map] (factor) -- (mword);
\draw[map] (eray) -- (witness);
\draw[map] (mword) -- (witness);
\draw[map] (witness) -- node[above,font=\small] {\(\mathscr O\)} (receiver);

\node[note,below=13mm of witness] (fibre) {For fixed \(p,r,s,\delta\), all allowed
\(t\) map to the same witness; retaining \(t\) recovers the source.};
\draw[map,dashed] (eray.south) -- (fibre.north west);
\node[note,below=13mm of receiver] (boundary) {No reverse arrow is proved:
receiver invertibility does not create the initial trace source.};
\draw[map,dashed] (receiver.south) -- (boundary.north);
\end{tikzpicture}
}
\caption{The typed composite from a marked mod-\(24\) rational trace to an
integral witness and then to its literal receiver.  The two arithmetic
branches are retained, as are the source fibre and the missing reverse map.}
\label{fig:mod24-receiver-composite}
\end{figure}

